\unit{Exposição 221 -- Os esquemas de Hilbert}
\sid{221}
\src{249}
\seminares{13.º ano, 1960/61, n.º 221\\Maio de 1961}
\paperhead{TÉCNICAS DE CONSTRUÇÃO E TEOREMAS DE EXISTÊNCIA EM GEOMETRIA ALGÉBRICA\\
IV. OS ESQUEMAS DE HILBERT}%
  {por Alexander GROTHENDIECK}

\subsection*{Introdução}
\addcontentsline{toc}{subsection}{Introdução}

As técnicas expostas em [2], I e II, eram essencialmente independentes de
qualquer hipótese de projetividade sobre os esquemas considerados.
Infelizmente, elas não permitem atualmente resolver os problemas de existência
formulados na Exposição II de [2]. Na presente exposição e na seguinte,
resolveremos esses problemas sob hipóteses de projetividade. As técnicas
empregadas são tipicamente projetivas e praticamente não utilizam os
resultados de I e II de [2]. Aqui construiremos os «esquemas de Hilbert»,
destinados a substituir o uso das coordenadas de Chow, como foi dito em [2],
II, n.º 2. Na exposição seguinte, a teoria da passagem ao quociente em
esquemas desenvolvida em III, juntamente com a teoria dos esquemas de Hilbert,
nos permitirá, por exemplo, construir os esquemas de Picard (definidos em [2],
II, n.º 3) sob condições bastante gerais.

Em resumo, pode-se dizer que agora se dispõe de uma técnica de construções
projetivas bastante satisfatória, salvo que ainda falta\footnote{Ver o adendo
ao fim da exposição.} um teorema de passagem ao quociente por grupos como o
grupo projetivo, atuando «sem pontos fixos» (cf. [2], III, n.º 8). A situação
parece até um pouco melhor em geometria analítica (quando nos limitamos ao
estudo dos espaços analíticos «projetivos» sobre um espaço analítico dado),
pois, para os espaços analíticos, desaparece a dificuldade da passagem ao
quociente por um grupo que atua de maneira adequada. Além disso, tanto em
geometria algébrica como em geometria analítica, restaria desenvolver uma
técnica de construção válida sem hipóteses de projetividade.

\fsection{1}{Famílias limitadas de feixes: propriedades de permanência}

Sejam $k$ um corpo e $X$ um $k$-pré-esquema, de tipo finito, para simplificar.
Para toda extensão $K/k$, obtém-se, portanto, um $K$-pré-esquema
$X_K=X\otimes_k K$. Se $\mathcal F$ é um feixe coerente em $X_K$ e $K'$ é
uma extensão de $K$, então $\mathcal F\otimes_K K'=\mathcal F_{K'}$ é um
feixe quase-coerente em $X_K\otimes_K K'=X_{K'}$. Posto isso, se $K$ e $K'$
são duas extensões quaisquer de $k$, $\mathcal F$ é um feixe quase-coerente em
$X_K$ e $\mathcal F'$
\src{250}
um feixe quase-coerente sobre $X_{K'}$, diremos que $\mathcal F$ e
$\mathcal F'$ são \emph{equivalentes} se existem $k$-homomorfismos de $K$,
$K'$ em uma mesma extensão $K''$ de $k$, tais que $\mathcal F_{K''}$ e
$\mathcal F'_{K''}$ sejam isomorfos sobre $X_{K''}$. Trata-se de uma relação de
equivalência, e interessar-nos-emos pelas classes de equivalência de feixes
para essa relação e pelos conjuntos de classes de equivalência. Observe-se que,
se $X_0$ é de tipo finito sobre $k$, então toda classe de feixes
coerentes pode ser definida por um feixe coerente sobre algum $X_K$, onde $K$
é uma extensão de tipo finito de $k$. Pode-se então, na definição das classes
de feixes coerentes, restringir-se às extensões algebricamente fechadas de $k$,
e pode-se também restringir a uma extensão algebricamente fechada fixa $\Omega$
de $k$, de grau de transcendência infinito, sendo dois feixes coerentes
$\mathcal F$ e $\mathcal F'$ sobre $X_\Omega$ equivalentes se e somente se
existe um $k$-automorfismo $\sigma$ de $\Omega$ tal que
$\mathcal F\otimes_\Omega(\Omega,\sigma)$ seja isomorfo a $\mathcal F'$.
Verifica-se que há uma correspondência biunívoca entre as classes de feixes
coerentes para uma ou outra definição.

Sejam $E$ e $E'$ dois conjuntos de classes de feixes coerentes sobre $X$.
Consideremos as classes de todos os feixes da forma
$\mathcal F\otimes\mathcal F'$, onde $\mathcal F$ e $\mathcal F'$ são feixes
coerentes sobre um mesmo $X_K$, com a classe de $\mathcal F$ em $E$ e a de
$\mathcal F'$ em $E'$. Obtém-se assim um conjunto de classes de feixes
coerentes, que será denotado por $E\otimes E'$. Definem-se do mesmo modo os
$\operatorname{Tor}_i(E,E')$, etc. De modo geral, a toda função $U$ que associa
a toda sequência $\mathcal F_1,\ldots,\mathcal F_n$ de $n$ feixes coerentes
sobre um mesmo $X_K$ um conjunto $U(\mathcal F_1,\ldots,\mathcal F_n)$ de
feixes coerentes sobre $X_K$, e que tem uma propriedade evidente de
compatibilidade com os isomorfismos de feixes e as imagens inversas por mudança
de base, associa-se uma função, denotada pelo mesmo símbolo $U$, que associa a
toda sequência de $n$ conjuntos $E_1,\ldots,E_n$ de classes de feixes
coerentes um conjunto $U(E_1,\ldots,E_n)$ de classes de feixes coerentes.

Nosso objetivo neste item é definir certos conjuntos de classes de feixes,
qualificados de \emph{limitados}, e mostrar que as operações $U$ mais usuais,
aplicadas a conjuntos limitados, ainda produzem conjuntos limitados.

Seja $X$ um pré-esquema de tipo finito sobre $S$ Noetheriano. Para todo
$s\in S$, a fibra $X_s$ é um pré-esquema de tipo finito sobre $k(s)$, e
consideraremos classes de feixes coerentes sobre $X_s$ no sentido precedente.
Isso dá sentido à expressão «classe de feixes coerentes sobre uma fibra de
$X/S$» e às expressões análogas. Do mesmo modo, procedendo separadamente sobre
cada fibra, podem-se ainda considerar operações tais como $E\otimes E'$, etc.,
que fazem corresponder a sistemas de conjuntos de
\src{251}
classes de feixes coerentes sobre as fibras de $X/S$, outro conjunto de classes
de feixes coerentes sobre as fibras de $X/S$.

\result{DEFINIÇÃO}{1.1} Seja $E$ um conjunto de classes de feixes coerentes
sobre as fibras de $X/S$. Diz-se que $E$ é \emph{limitado} se existem um
pré-esquema $S'$ de tipo finito sobre $S$ e um feixe coerente $\mathcal F'$ sobre
$X'=X\times_S S'$ tais que $E$ esteja contido no conjunto de classes de feixes
sobre as fibras de $X/S$ definido por $\mathcal F'$.

Por definição, este último faz corresponder a cada $s\in S$ as classes dos
feixes $\mathcal F'\otimes_{S'}k(s')$, em que $s'$ percorre os pontos de $S'$
situados sobre $s$ (de modo que $k(s')$ é uma extensão de $k(s)$ e
$X_s\otimes_{k(s)}k(s')$ identifica-se com a fibra
$X'\otimes_{S'}k(s')=X'_{s'}$ de $X'$ em $s'$). Pode-se dizer que as famílias
limitadas são aquelas contidas em uma \emph{família algébrica} de feixes
coerentes parametrizada por um $S'$ de tipo finito sobre $S$.

Uma união finita de famílias limitadas é limitada (tome-se como pré-esquema a
soma disjunta dos pré-esquemas de parâmetros $S_i$ que definem as famílias
algébricas majorantes). Uma mudança de base $T\longrightarrow S$ transforma uma
família limitada relativa a $X/S$ em uma família limitada relativa a $X_T/T$, e
a recíproca é verdadeira se $T\longrightarrow S$ é sobrejetivo (ou, mais
geralmente, se sua imagem contém os $s$ que intervêm efetivamente na família
dada $E$ para $X/S$). Isso reduz, em princípio, a determinação das famílias
limitadas ao caso em que $S$ é o espectro de uma álgebra de tipo finito sobre o
anel dos inteiros $\mathbf Z$.

Se $E$ e $E'$ são famílias limitadas de classes de feixes relativas a $X/S$,
então $E\otimes E'$ também é limitada. Com efeito, se $E$ e $E'$ são majoradas,
respectivamente, pelas famílias algébricas definidas por $T\longrightarrow S$ e
$\mathcal F$ sobre $X_T$, e por $T'\longrightarrow S$ e $\mathcal F'$ sobre
$X_{T'}$, vê-se que $E\otimes E'$ é majorada pela família algébrica definida por
$T''=T\times_S T'\longrightarrow S$ e pelo feixe $\mathcal F''$ sobre
$X_{T''}$, produto tensorial das imagens inversas de $\mathcal F$ e
$\mathcal F'$ sobre $X_{T''}$. Esse raciocínio só é válido porque o funtor
$\mathcal F\otimes\mathcal F'$ nas variáveis $\mathcal F$, $\mathcal F'$ é exato
à direita e, portanto, comuta com a extensão da base (e, em particular, com a
passagem às fibras). Ele não se aplica, tal como está, às operações locais como
$\operatorname{Tor}_i(E,E')$, $\operatorname{Hom}(E,E')$ e
$\operatorname{Ext}^i(E,E')$. Pode-se, contudo, mostrar que essas operações
continuam transformando conjuntos limitados em conjuntos limitados, procedendo
como para $E\otimes E'$, mas utilizando também resultados do tipo seguinte
(todos contidos em [3], IV, 6.11): uma família limitada $E$ é sempre majorada
por uma família algébrica definida por um feixe coerente $\mathcal F$ sobre um
$X_T$ ($T$ de tipo finito sobre $S$) que é plano relativamente a $T$. (O espaço
de parâmetros inicial é «cortado em partes».) Tais
\src{252}
propriedades de planicidade em feixes adequados asseguram, com efeito, a
comutação de operações como $\operatorname{Tor}_i(\mathcal F,\mathcal F')$
com uma mudança de base arbitrária. O mesmo método aplica-se a operações de
natureza global: imagens diretas e imagens diretas derivadas de feixes
coerentes por morfismos próprios, $\operatorname{Ext}$ globais relativos a
morfismos próprios (cf. [5], III, \S~6), etc.; todas essas operações
transformam famílias limitadas de feixes em famílias limitadas de feixes.
(N. B. -- Aqui os pré-esquemas sobre os quais se tomam os diversos feixes
podem mudar pelas operações do tipo considerado.)

Os dois enunciados seguintes demonstram-se essencialmente pela mesma técnica
de planicidade; para a decomposição primária nas fibras de um morfismo de tipo
finito, veja-se em particular [5], IV.

\result{PROPOSIÇÃO}{1.2} Sejam $E$ e $E'$ conjuntos limitados de classes de
feixes nas fibras de $X/S$, supondo-se $X$ próprio sobre $S$. Então:

\medskip
\noindent (i) A família dos núcleos, conúcleos e imagens de homomorfismos
$\mathcal F\longrightarrow\mathcal F'$, em que a classe de $\mathcal F$
pertence a $E$ e a de $\mathcal F'$ pertence a $E'$, é limitada.

\smallskip
\noindent (ii) A família dos feixes $\mathcal F''$ que são extensões de um
$\mathcal F$ por um $\mathcal F'$, em que a classe de $\mathcal F$ pertence
a $E$ e a de $\mathcal F'$ pertence a $E'$, é limitada.

Mediante uma mudança de base adequada, pode-se supor que $E$ e $E'$ sejam
definidos, respectivamente, por um feixe coerente $\mathcal G$ e um feixe
coerente $\mathcal G'$ sobre um $X_T/T$, sendo $T$ de tipo finito sobre $S$.
Além disso, podem-se supor satisfeitas certas hipóteses de planicidade, que
implicam que a formação dos feixes
$f_{T*}(\operatorname{Hom}_{X_T}(\mathcal G,\mathcal G'))$ e
$\operatorname{Ext}^{1}_{f_T}(\mathcal G,\mathcal G')$ comuta com a mudança
de base por um morfismo arbitrário $T'\longrightarrow T$. Pode-se também
supor que os feixes coerentes precedentes sobre $T$ sejam localmente livres.
Sejam então $T_0$ e $T_1$ os fibrados vetoriais sobre $T$ cujos feixes de
germes de seções são, respectivamente, os feixes precedentes. Define-se então
canonicamente um homomorfismo
$\mathcal G_{T_0}\longrightarrow\mathcal G'_{T_0}$ de feixes coerentes sobre
$X_{T_0}$ e uma extensão
\[
0\longrightarrow\mathcal G_{T_1}\longrightarrow\mathcal G''
\longrightarrow\mathcal G'_{T_1}\longrightarrow0
\]
de feixes coerentes sobre $X_{T_1}$, dotada de uma propriedade universal
evidente. Esse segundo feixe define uma família algébrica que majora a família
considerada em (ii). O mesmo vale para o núcleo, o conúcleo e a imagem do
homomorfismo
\src{253}
precedente, cuja consideração estabelece (i) (desde que se suponha o conúcleo
plano relativamente a $T_0$, caso ao qual ainda se reduz decompondo $T_0$ em
partes\ldots).

\result{PROPOSIÇÃO}{1.3} Seja $E$ uma família limitada de classes de feixes
nas fibras de $X/S$. Então as classes dos feixes estruturais dos
$(\operatorname{supp}\mathcal F)_{\mathrm{red}}$, em que $\mathcal F$ é um
feixe coerente sobre um $X_K$, com $K$ algebricamente fechado, cuja classe
pertence a $E$, formam uma família limitada.

Aqui, $(\operatorname{supp}\mathcal F)_{\mathrm{red}}$ designa o suporte de
$\mathcal F$, munido da estrutura reduzida induzida, isto é, seu feixe
estrutural é o quociente de $\mathcal O_{X_K}$ pelo maior feixe de ideais que
define $\operatorname{supp}\mathcal F$. Pode-se demonstrar o resultado
análogo a (1.3) para os feixes deduzidos canonicamente de $\mathcal F$ pela
teoria da decomposição primária; por exemplo, os
$\mathcal F/\mathcal F_i$, em que os $\mathcal F_i$ são os subfeixes de
$\mathcal F$ primários para as componentes do suporte de $\mathcal F$ e
mínimos com essa propriedade; ou os $\mathcal O_{X_K}/\mathfrak p$, em que
$\mathfrak p$ é um feixe de ideais primos associado a $\mathcal F$; ou os
$\mathcal O_{X_K}/\mathfrak q$, em que $\mathfrak q$ é um feixe de ideais
primários associado a uma componente do suporte de $\mathcal F$ (sendo o
corpo de referência algebricamente fechado).

\fsection{2}{Famílias limitadas e polinômio de Hilbert}

No que segue, supomos que $X$ seja projetivo sobre $S$ e esteja munido de um
feixe muito amplo, denotado por $\mathcal O_X(1)$. Para toda extensão $K$ do
corpo residual $k(s)$ de um ponto $s$ de $S$, considera-se sobre $X_K$ o feixe
correspondente $\mathcal O_{X_K}(1)$, que será ainda muito amplo.

A todo feixe coerente $\mathcal F$ sobre $X_K$ associa-se a função
\[
P_{\mathcal F}(n)=\text{característica de Euler--Poincaré de $\mathcal F(n)$
sobre $X_K$},
\]
que é um polinômio no inteiro $n$, chamado \emph{polinômio de Hilbert} de
$\mathcal F$. Para valores grandes de $n$, $P(n)$ nada mais é do que a
dimensão de $H^0(X_K,\mathcal F(n))$ sobre $K$, pois os
$H^i(X_K,\mathcal F(n))$ são nulos para $i>0$ e $n$ grande.

Se agora $\mathcal F$ for um feixe coerente sobre $X$, plano relativamente a
$S$, então os polinômios de Hilbert dos feixes $\mathcal F_s$ induzidos nas
fibras $X_s$ relativas a uma mesma componente conexa de $S$ serão todos iguais
[5], III, \S~7. Daí resulta (sem hipótese de planicidade) que o conjunto dos
polinômios de Hilbert dos feixes $\mathcal F_s$, $s\in S$, é finito para
todo feixe coerente $\mathcal F$ sobre $X$.
\src{254}
Recordemos, por outro lado, que, se $\mathcal F$ é um feixe coerente sobre $X$,
este é isomorfo a um feixe quociente de um feixe $\mathcal O_X(-n)^N$, para
$n$ e $N$ suficientemente grandes. Portanto, os feixes $\mathcal F_s$
induzidos nas fibras são igualmente quocientes do feixe $\mathcal O(-n)$ na
fibra.

Das duas observações, deduz-se a parte «somente se» do teorema seguinte:

\result{TEOREMA}{2.1} Seja $X$ projetivo sobre um $S$ Noetheriano, e seja
$\mathcal O_X(1)$ muito amplo sobre $X$ relativamente a $S$. Seja $E$ um
conjunto de classes de feixes nas fibras de $X/S$. Para que $E$ seja limitado,
é necessário que
satisfaça as condições seguintes:

\medskip
\noindent (a) Existe um feixe coerente $\mathcal L$ sobre $X$ (que se pode
supor da forma $\mathcal O_X(-n)^N$) tal que $E$ esteja contido na família das
classes de feixes coerentes quocientes de feixes da forma $\mathcal L_K$.

\smallskip
\noindent (b) Os polinômios de Hilbert $P_{\mathcal F}$ dos feixes
$\mathcal F$ cuja classe pertence a $E$ são elementos de um mesmo conjunto
finito de polinômios.

Resta provar a parte «se», que será um caso particular de um resultado mais
preciso. Para todo Módulo coerente $\mathcal F$ sobre um pré-esquema de tipo
finito sobre um corpo $K$, e todo inteiro $r$, seja $N_r$ o sub-Módulo de
$\mathcal F$ cujas seções sobre um aberto são as seções de $\mathcal F$ sobre
esse aberto cujo suporte tem dimensão $<r$. Tem-se, portanto,
$N_r=\mathcal F$ para $r>\dim\operatorname{supp}\mathcal F$, $N_r=0$ para
$r\leq0$, e obtém-se uma filtração crescente finita de $\mathcal F$ cujos
fatores $N_r/N_{r+1}$ têm como ciclos primos associados exatamente os ciclos primos
associados a $\mathcal F$ que têm dimensão $r$. Poremos
\[
\mathcal F_{(r)}=\mathcal F/N_r.
\]
Assim, $\mathcal F_{(r)}$ tem como ciclos primos associados exatamente os
ciclos primos associados a $\mathcal F$ que têm dimensão $\geq r$ e, em
particular, é igual a $\mathcal F$ se e somente se os ciclos primos associados
a $\mathcal F$ têm dimensão $\geq r$. Posto isso:

\result{TEOREMA}{2.2} Nas condições de (2.1), seja $s$ um inteiro e suponhamos
que $E$ satisfaça a condição (a) e a seguinte forma enfraquecida de (b):

\medskip
\noindent $(b_s)$ Os polinômios de Poincaré $P_{\mathcal F}$ dos feixes
$\mathcal F$ cuja classe pertence a $E$ têm coeficientes em graus
$\leq s-1$ que permanecem limitados.

Nessas condições, os feixes $\mathcal F_{(s)}$ (quando a classe de
$\mathcal F$ percorre $E$) formam uma família limitada. Além disso, os
coeficientes de grau $s-2$ dos $P_{\mathcal F}$ permanecem limitados
inferiormente.
\src{255}
Assim:

\result{COROLÁRIO}{2.3} Suponhamos que os feixes $\mathcal F$ cuja classe
pertence a $E$ sejam tais que todos os ciclos primos associados tenham dimensão
$d$ com $s\leq d\leq r$. Então, em (2.1) (b), pode-se considerar apenas os
coeficientes dos $P_{\mathcal F}$ em graus compreendidos entre $s-1$ e $r$.

O fim deste item é dedicado ao esboço da demonstração de (2.2). Os lemas-chave
são os dois lemas seguintes, o primeiro dos quais é bem conhecido (e resume o
conteúdo matemático útil das coordenadas de Chow):

\result{LEMA 2.4 (CHOW)}{} Consideremos os feixes estruturais de subesquemas
$Y$ de fibras $X_K$ ($K$ percorrendo extensões algebricamente fechadas de
corpos residuais de $S$), com $Y$ reduzido e todas as suas componentes de mesma
dimensão $r$ (considerando-se $\mathcal O_Y$ como um feixe quociente de
$\mathcal O_X$). Se os graus dos $Y$ permanecem limitados, os $Y$ formam uma
família limitada.

Aqui, o grau $a$ de $Y$ pode ser definido mais convenientemente pelo
coeficiente do termo dominante de
\[
P_{\mathcal O_Y}=a\frac{n^r}{r!}+\cdots.
\]

\result{LEMA}{2.5} Sejam $\mathcal L$ um feixe coerente sobre $X$ e $E$ um
conjunto de classes de feixes quocientes $\mathcal F$ de feixes $\mathcal L_K$
($K$ uma extensão de corpo residual de $S$). Suponhamos que as fibras de $X$
sobre $S$ tenham dimensão $\leq r$, e ponhamos
\[
P_{\mathcal F}(n)=a_{\mathcal F}\frac{n^r}{r!}
+b_{\mathcal F}\frac{n^{r-1}}{(r-1)!}
+\text{termos de grau $<r-1$}.
\]
Então o coeficiente $a_{\mathcal F}$ permanece limitado, e $b_{\mathcal F}$
permanece limitado inferiormente. Se $b_{\mathcal F}$ permanece limitado,
então a família dos $\mathcal F_{(r)}$ é limitada.

\emph{Demonstração.} -- Substituindo, se necessário, $S$ por uma união de
subesquemas de $S$ que recobre $S$, pode-se supor que existe um morfismo finito
$f:X\longrightarrow\mathbf P^r_S$, tal que $\mathcal O_X(1)$ seja isomorfo à
imagem inversa de $\mathcal O_{\mathbf P^r_S}(1)$; portanto, para todo feixe
coerente $\mathcal F$ sobre $X$, tem-se
$P_{\mathcal F}=P_{f_*(\mathcal F)}$. Por outro lado, verifica-se facilmente
(pela técnica do item precedente) que um conjunto de feixes $\mathcal F$ sobre
$X$ é limitado se e somente se o conjunto dos $f_*(\mathcal F)$ o é. Por fim,
\[
f_*(\mathcal F)_{(r)}=f_*(\mathcal F_{(r)}).
\]
\src{256}
Isso nos reduz, portanto, ao caso $X=\mathbf P^r_S$. Além disso, pode-se supor
que $\mathcal L=\mathcal O_{\mathbf P^r_S}(k)^s$ para $k$, $s$ convenientes. O
coeficiente $a_{\mathcal F}$ satisfaz
\[
0\leq a_{\mathcal F}\leq s
\]
e, portanto, permanece limitado. Dito isso, é equivalente afirmar que os
$P_{\mathcal F}(n)$ têm um coeficiente de $n^{r-1}$ que permanece limitado
inferiormente (resp. limitado), ou afirmar o mesmo para
$P_{\mathcal F}(n-k)=P_{\mathcal F(-k)}(n)$. Isso nos reduz ao caso em que
\[
\mathcal L=\mathcal O_{\mathbf P^r_S}^{s}.
\]

Consideremos a sequência exata
\[
0\longrightarrow N_r\longrightarrow\mathcal F
\longrightarrow\mathcal F_{(r)}\longrightarrow0,
\]
da qual se obtém
\[
P_{\mathcal F}=P_{\mathcal F_{(r)}}+P_{N_r},
\]
e, como o coeficiente de $n^{r-1}$ em $P_{N_r}$ é não negativo (pois
$\dim\operatorname{supp}N_r\leq r-1$), tem-se
\[
b_{\mathcal F_{(r)}}\leq b_{\mathcal F}.
\]
Isso nos permite, para demonstrar as afirmações do lema, substituir
$\mathcal F$ por $\mathcal F_{(r)}$, isto é, supor que os feixes quocientes
$\mathcal F$ de $\mathcal L$ considerados são livres de torção.

Como $\mathbf P^r_K$ é normal, segue-se que $\mathcal F$ é localmente livre de
posto $a=a_{\mathcal F}$ em um aberto $U=\mathbf P^r_K-Y$, onde $Y$ tem
codimensão $\geq2$. Logo, $\bigwedge^a\mathcal F$ é um feixe sobre
$\mathbf P^r_K$ cuja restrição a $U$ é invertível; portanto (como
$\mathbf P^r_K$ é regular e $Y$ tem codimensão $\geq2$), essa restrição é
isomorfa à restrição de um feixe invertível sobre $\mathbf P^r_K$, determinado
salvo isomorfismo. Este último é da forma $\mathcal O_{\mathbf P^r_K}(d)$ para um
inteiro $d$ bem determinado. Como $\bigwedge^a\mathcal F$ é um quociente de
$\bigwedge^a\mathcal O_{\mathbf P^r_K}^{n}\simeq
\mathcal O_{\mathbf P^r_K}^{N}$, com $N=\binom{s}{a}$, possui $N$ seções canônicas,
que definem, portanto, seções de
\src{257}
$\mathcal O_{\mathbf P^r_K}(d)$ sobre $U$, que são restrições de seções
$s_i$ ($1\leq i\leq N$) de $\mathcal O_{\mathbf P^r_K}(d)$ (pois
$\mathbf P^r_K$ é normal e $Y$ tem codimensão $\geq2$). Estas seções
geram $\mathcal O_{\mathbf P^r_K}(d)$ nos pontos de $U$, portanto não são
todas nulas, o que implica que $d\geq0$. Por outro lado, um cálculo simples
mostra que
\[
b_{\mathcal F}=a_{\mathcal F}(r+1)/2+d.
\]
Isso mostra, em particular, que $b_{\mathcal F}>0$, portanto
$b_{\mathcal F}$ é limitado inferiormente. Esse coeficiente permanece limitado
se e somente se $d$ permanece limitado; mostremos que, nesse caso, $\mathcal F$
pertence a uma família limitada. Podemos supor que $a_{\mathcal F}$ e
$b_{\mathcal F}$ sejam fixos, digamos iguais a $a$ e $b$; portanto, $d$ é fixo.
O conhecimento das $N$ seções
$s_i$ de $\mathcal O_{\mathbf P^r_K}(d)$, isto é, de um homomorfismo
\[
s:\bigwedge^a\mathcal L_K\longrightarrow\mathcal O_{\mathbf P^r_K}(d),
\]
permite recuperar $\mathcal F$ como a coimagem do homomorfismo composto
correspondente:
\[
\mathcal L_K\longrightarrow
\operatorname{Hom}\!\left(\bigwedge^{a-1}\mathcal L_K,
\bigwedge^a\mathcal L_K\right)\longrightarrow
\operatorname{Hom}\!\left(\bigwedge^{a-1}\mathcal L_K,
\mathcal O_{\mathbf P^r_K}(d)\right),
\]
em que o primeiro é o homomorfismo canônico proveniente do produto exterior, e
o segundo é deduzido de $s$. A conclusão decorre então de (1.2) (i).

Combinando os dois lemas precedentes, obtém-se:

\result{LEMA}{2.6} Sob as hipóteses preliminares de (2.1), suponhamos que, para
todo $\mathcal F$, se tenha
\[
P_{\mathcal F}(n)=a_{\mathcal F}\frac{n^r}{r!}
+b_{\mathcal F}\frac{n^{r-1}}{(r-1)!}
+\text{termos de grau $<r-1$},
\]
e que os coeficientes $a_{\mathcal F}$ permaneçam limitados. Então os
coeficientes $b_{\mathcal F}$ permanecem limitados inferiormente. Se os
$b_{\mathcal F}$ forem limitados, então os $\mathcal F_{(r)}$ formam uma família
limitada.

Podemos supor que os corpos de base $K$ dos feixes $\mathcal F$ sejam
algebricamente fechados. Munamos cada
$\operatorname{supp}\mathcal F_{(r)}$, união das componentes de grau
$r$ de
$\operatorname{supp}\mathcal F$, da estrutura reduzida induzida. Então os graus dos
$\operatorname{supp}\mathcal F_{(r)}$ permanecem limitados superiormente por
$a$; portanto, em virtude de (2.4), os
$\operatorname{supp}\mathcal F_{(r)}$ formam um conjunto limitado. Além disso,
para toda componente de $\operatorname{supp}\mathcal F_{(r)}$, o comprimento de
$\mathcal F_{(r)}$ nessa componente é $\leq a$; portanto, se
$\mathcal J_{\mathcal F}$ é o Ideal que define
$\operatorname{supp}\mathcal F_{(r)}$, então $\mathcal F_{(r)}$
\src{258}
pode ser considerado um Módulo sobre o subesquema $Y_{\mathcal F}$ de $X$
definido por $\mathcal J_{\mathcal F}^{a}$. Como no lema precedente, reduzimo-nos também
ao caso em que $\mathcal F=\mathcal F_{(r)}$, de modo que $\mathcal F$
provém de um Módulo sobre $Y_{\mathcal F}$. Os $Y_{\mathcal F}$ correspondem
a uma família limitada de Módulos quocientes dos $\mathcal O_{X_K}$ e,
portanto, provêm de um subesquema fechado $Y$ de um esquema $X\times_S T$.
Pode-se então aplicar (2.5) a $Y/T$ e $\mathcal L\otimes_XY$, o que dá a
conclusão.

Podemos agora demonstrar (2.2) por indução sobre o limite superior $r$ dos
$\dim\operatorname{supp}\mathcal F$. O enunciado é trivial para $r<0$;
suponhamos, pois, que $r\geq0$ e que o enunciado tenha sido demonstrado para
$r'<r$. Em virtude de (2.6), os $\mathcal F_{(r)}$ formam um conjunto
limitado; portanto, também em virtude de (1.2) (i), os núcleos dos
homomorfismos $\mathcal L_K\longrightarrow\mathcal F_{(r)}$ formam um
conjunto limitado. Existe, pois, um Módulo coerente $\mathcal L'$ sobre $X$
tal que os núcleos em questão e, consequentemente, também os
\[
N_r(\mathcal F)=\operatorname{Ker}(\mathcal F\longrightarrow\mathcal F_{(r)}),
\]
sejam quocientes dos Módulos $\mathcal L'_K$. Como os $\mathcal F_{(r)}$ são
limitados, os $P_{\mathcal F_{(r)}}$ permanecem limitados, e a fórmula
\[
P_{\mathcal F}=P_{\mathcal F_{(r)}}+P_{N_r}
\]
mostra então que os $P_{N_r}$ satisfazem a mesma condição $(b_s)$ que os
$P_{\mathcal F}$. Logo, os $N_r$ satisfazem as condições (a) e $(b_s)$ e,
pela hipótese de indução, os $(N_r)_{(s)}$ permanecem limitados. Ora,
$\mathcal F_{(s)}$ é uma extensão de $\mathcal F_{(r)}$ por
$(N_r)_{(s)}$; portanto, por (1.2) (ii), os $\mathcal F_{(s)}$ permanecem
limitados. Para a última afirmação de (2.2), observa-se que os núcleos $N_s$ de
$\mathcal F\longrightarrow\mathcal F_{(s)}$ permanecem limitados em virtude
de (1.2) (i), e que o coeficiente do termo em $n^{s-1}$ em $P_{N_s}$ permanece
limitado; o lema 2.6 demonstra então que o coeficiente do termo seguinte
permanece minorado. Isso conclui a demonstração.

\fsection{3}{Esquemas de Hilbert: definição, teorema de existência}

Sejam $X$ um pré-esquema sobre outro $S$, e $\mathcal F$ um Módulo
quase-coerente sobre $X$. Denotemos por
\[
\operatorname{Quot}(\mathcal F/X/S)
\]
o conjunto dos Módulos quase-coerentes quocientes de $\mathcal F$ que são
planos sobre $S$. Seja agora $S'\longrightarrow S$ um morfismo de mudança de
base; ponhamos $X'=X\times_S S'$ e
$\mathcal F'=\mathcal F\otimes_{\mathcal O_S}\mathcal O_{S'}$; assim, $X'$
é um pré-esquema sobre $S'$ munido de um Módulo
\src{259}
quase-coerente $\mathcal F'$, e pode-se considerar
$\operatorname{Quot}(\mathcal F'/X'/S')$. Ponhamos
\[
\underline{\operatorname{Quot}}_{\mathcal F/X/S}(S')
=\operatorname{Quot}(\mathcal F'/X'/S')
\qquad\text{(onde $X'=X\times_SS'$)}.
\]
Se agora $S''\longrightarrow S'$ for um $S$-morfismo, então
$X''=X\times_SS''$ é isomorfo a $X'\times_{S'}S''$, e $\mathcal F''$ é
isomorfo a $\mathcal F'\otimes_{\mathcal O_{S'}}\mathcal O_{S''}$. Como o
funtor imagem inversa
\[
\mathcal G'\longmapsto
\mathcal G'\otimes_{\mathcal O_{S'}}\mathcal O_{S''}
\]
da categoria dos Módulos quase-coerentes sobre $X'$ para a categoria dos
Módulos quase-coerentes sobre $X''$ é exato à direita e transforma Módulos
$S'$-planos em Módulos $S''$-planos, obtém-se uma aplicação natural
\[
\operatorname{Quot}(\mathcal F'/X'/S')\longrightarrow
\operatorname{Quot}(\mathcal F''/X''/S''),
\]
de modo que $\underline{\operatorname{Quot}}_{\mathcal F/X/S}(S')$ é um
funtor contravariante em $S'$ (pré-esquema sobre $S$), com valores na categoria
dos conjuntos. Doravante, suporemos que $X$ seja projetivo sobre um $S$
noetheriano, que $\mathcal F$ seja coerente e, para simplificar, limitar-nos-emos
aos $S'$ sobre $S$ que sejam localmente noetherianos.

\result{TEOREMA}{3.1} Nessas condições, o funtor contravariante
$\underline{\operatorname{Quot}}_{\mathcal F/X/S}$ sobre a categoria dos
$S$-pré-esquemas localmente noetherianos é representável por um
$S$-pré-esquema $\operatorname{Quot}_{\mathcal F/X/S}$, união disjunta de uma
sequência de $S$-esquemas projetivos (a fortiori,
$\operatorname{Quot}_{\mathcal F/X/S}$ é localmente de tipo finito sobre
$S$).

Obteremos tal decomposição da seguinte maneira. Seja $\mathcal O_X(1)$ um
feixe invertível sobre $X$, muito amplo relativamente a $S$. Para todo
polinômio $P(n)$ com coeficientes racionais, seja
$\operatorname{Quot}^P(\mathcal F/X/S)$ a parte de
$\operatorname{Quot}(\mathcal F/X/S)$ formada pelos quocientes coerentes
$\mathcal G$ de $\mathcal F$ que são planos sobre $S$ e cujo polinômio de
Hilbert em todo $s\in S$ é igual a $P$. Poremos então
\[
\underline{\operatorname{Quot}}^P_{\mathcal F/X/S}(S')
=\operatorname{Quot}^P(\mathcal F'/X'/S')
\]
e obteremos assim um subfuntor de
$\underline{\operatorname{Quot}}_{\mathcal F/X/S}$. A propriedade de
invariância dos polinômios de Hilbert recordada no número 2 implica o seguinte:
para que $\underline{\operatorname{Quot}}_{\mathcal F/X/S}$ seja
representável, é necessário e suficiente que os
$\underline{\operatorname{Quot}}^P_{\mathcal F/X/S}$ o sejam; nesse caso, o
$S$-pré-esquema $\operatorname{Quot}_{\mathcal F/X/S}$ que o representa é
isomorfo ao pré-esquema união disjunta dos
$\operatorname{Quot}^P_{\mathcal F/X/S}$ que representam os funtores
$\underline{\operatorname{Quot}}^P_{\mathcal F/X/S}$. Posto isso, (3.1)
será, portanto, consequência do teorema seguinte:
\src{260}
\result{TEOREMA}{3.2} Com as notações anteriores, o funtor
$\underline{\operatorname{Quot}}^P_{\mathcal F/X/S}$ é representável por um
$S$-pré-esquema projetivo $\operatorname{Quot}^P_{\mathcal F/X/S}$.

O restante deste número é dedicado à demonstração de (3.2).

Seja $\nu$ um inteiro. Para todo $S'$ sobre $S$, denotamos por $A_\nu(S')$ o
conjunto dos quocientes $\mathcal G=\mathcal F'/\mathcal K$ de
$\mathcal F'=\mathcal F\otimes_{\mathcal O_S}\mathcal O_{S'}$, que são
coerentes, planos sobre $S'$ e satisfazem as condições seguintes:
\begin{align*}
\text{a.}\quad &R^if'_*(\mathcal G(n))=0 &&\text{para $i>0$ e $n\geq\nu$},\\
\text{b.}\quad &R^if'_*(\mathcal K(n))=0 &&\text{para $i>0$ e $n\geq\nu$},\\
\text{c.}\quad &f'_*(\mathcal K(\nu+k))
=\mathcal S'_k f'_*(\mathcal K(\nu)) &&\text{para $k\geq0$}.
\end{align*}
Para esta última relação, supõe-se que $X$ esteja escrito como o espectro primo
homogêneo de uma álgebra graduada quase-coerente $\mathcal S_*$ sobre $S$,
graduada em graus positivos e gerada por $\mathcal S_1$, e põe-se
$\mathcal S'_* =\mathcal S_*\otimes_{\mathcal O_S}\mathcal O_{S'}$, de modo
que $X'$ seja o espectro primo homogêneo de $\mathcal S'_*$. Para demonstrar
(3.2), reduz-se facilmente ao caso em que $X=\mathbf P^r_S$ (pois $S$ é união
de abertos $U$ tais que $X|U$ seja um subesquema fechado de $\mathbf P^r_U$ e
$\mathcal O_X(1)$ seja induzido por $\mathcal O_{\mathbf P^r_U}(1)$), e ao
caso em que $\mathcal F$ é da forma
$\mathcal O_{\mathbf P^r_S}^{N}$, sendo, portanto, plano sobre $S$. Então, no
que precede, os feixes $\mathcal K$ também são planos sobre $S'$. Resulta das
relações de Künneth ([5], III, \S~7) que as condições (a) e (b) são estáveis
por mudança de base e implicam que, para $n\geq\nu$, a formação de
$f'_*(\mathcal G(n))$ e de $f'_*(\mathcal K(n))$ comute com a mudança de base
(loc. cit.); portanto, (c) também é estável por mudança de base.

Em outras palavras, $A_\nu(S')$ é um funtor contravariante em $S'$, mais
precisamente um subfuntor de
$A(S')=\underline{\operatorname{Quot}}^P_{\mathcal F/X/S}(S')$. Ao variar
$\nu$, obtém-se assim uma sequência crescente de subconjuntos $A_\nu(S')$ de
$A(S')$, cuja união é $A(S')$ em virtude de um teorema bem conhecido de
Serre [5], III, \S~2. Observemos agora que, se $\mathcal G$ é um quociente
coerente de $\mathcal F'$, plano sobre $S'$, e $s$ é um ponto de $S'$ tal que
a mudança de base $\operatorname{Spec}(k(s))\longrightarrow S'$ dá origem a
um quociente $\mathcal G_s$ de $\mathcal F_s$ que satisfaz as condições (a),
(b) e (c), isto é, que pertence a $A_\nu(\operatorname{Spec}(k(s)))$, então
existe uma vizinhança aberta $U$ de $s$ tal que essas mesmas condições sejam
verificadas para $\mathcal G|f'^{-1}(U)$, isto é, esse quociente pertence a
$A_\nu(U)$; para (a) e (b), isso resulta, com efeito, do «teorema das funções
holomorfas» [5], III, \S~7, e (c) resulta do lema de Nakayama e do fato de que,
em todo caso, sabe-se que
$f'_*(\mathcal K(n+k))=\mathcal S'_k f'_*(\mathcal K(n))$
\src{261}
para $n$ suficientemente grande e $k\geq0$ (cf. [5], III, \S~2). Dessas
observações conclui-se o seguinte (compare-se [4], IV). Para que o funtor $A$
seja representável, é necessário e suficiente que os funtores $A_\nu$ o sejam;
nesse caso, o $S$-pré-esquema $Q$ que representa $A$ é uma união crescente de
subpré-esquemas abertos $Q_\nu$ que representam os $A_\nu$.

Seja
\[
M_* =\bigoplus_{n\geq0}f_*(\mathcal F(n))=\mathcal S_*^{N},
\]
de modo que
\[
M'_* =M_*\otimes_{\mathcal O_S}\mathcal O_{S'}
=\bigoplus_{n\geq0}f'_*(\mathcal F'(n))=\mathcal S'_*{}^{N}.
\]
De (a) resulta que:

\noindent (a') $f'_*(\mathcal G(n))$ é localmente livre de posto $P(n)$ para
$n\geq\nu$, [5], III, \S~7, e, de (b), para $i=1$, que:

\noindent (a'') $f'_*(\mathcal G(n))$ é um Módulo quociente de $M'_n$.

Além disso, o conhecimento desse Módulo quociente, para $n=\nu$, determina,
em virtude de (c), os sub-Módulos $f'_*(\mathcal K(n))$ de $M'_n$ para
$n\geq\nu$ e, portanto, determina $\mathcal K$ e, por conseguinte, $\mathcal G$.
Obtém-se assim uma aplicação injetiva
\[
A_\nu(S')\longrightarrow\operatorname{Grass}_{P(\nu)}(M'_\nu)
\]
de $A_\nu(S')$ no conjunto dos Módulos quocientes localmente livres de posto
$P(\nu)$ de $M'_\nu$ e, portanto, um homomorfismo funtorial
\[
i_\nu:A_\nu(S')\longrightarrow
\underline{\operatorname{Grass}}_{P(\nu)}(M_\nu)(S'),
\]
onde o funtor do segundo membro é representado pelo esquema de Grassmann
$\operatorname{Grass}_{P(\nu)}(M_\nu)$ (compare-se [4], V), que é projetivo
sobre $S$. Afirmo:

\result{LEMA}{3.3} $A_\nu(S')$ é um funtor representável, e o morfismo
$Q_\nu\longrightarrow\operatorname{Grass}_{P(\nu)}(M_\nu)$ que representa o
homomorfismo $i_\nu$ é uma imersão (o que implica que $Q_\nu$ é
quase-projetivo sobre $S$).

Essa afirmação é equivalente à seguinte (compare-se [4], IV): suponha-se dado
um Módulo quociente $N$ de $M'_\nu$, localmente livre de posto $P(\nu)$; então
existe um subpré-esquema $Z$ de $S'$ tal que, para todo pré-esquema localmente
noetheriano $T'$
\src{262}
sobre $S'$, a imagem inversa de $N$ sobre $T'$ pertence a
$\operatorname{Im}A_\nu(T')$ se e somente se $T'\longrightarrow S'$ se
fatoriza pelo subpré-esquema $Z$. Mudando as notações, pode-se supor $S'=S$,
isto é, tem-se um quociente $N_\nu$ de $M_\nu$ por um sub-Módulo $R_\nu$.
Para que esse quociente provenha de um elemento de $A_\nu(S)$, é necessário e
suficiente que satisfaça as duas condições seguintes:

\medskip
\noindent (i) $M_{\nu+k}/\mathcal S_kR_\nu$ é localmente livre de posto
$P(\nu+k)$ para $k\geq0$.

\smallskip
\noindent (ii) O subfeixe $\mathcal K$ de $\mathcal F$ definido pelo
sub-Módulo graduado
\[
R_* =\sum_{k\geq0}\mathcal S_kR_\nu
\]
de $M_*$ ([5], II, \S~3) e o quociente $\mathcal G=\mathcal F/\mathcal K$
satisfazem as condições (a) e (b) acima.

Essas condições são manifestamente necessárias; por outro lado, se forem
satisfeitas, então o feixe $\mathcal G$ definido em (ii), uma vez que é
isomorfo ao feixe associado ao $\mathcal S_*$-Módulo graduado $N_*$ que é a
soma dos $M_{\nu+k}/\mathcal S_kR_\nu$, é plano sobre $S$ (pois suas fibras
são somandos diretos de localizações de $N$ em ideais primos homogêneos de
$\mathcal S_*$), e tem polinômio de Hilbert $P$ em virtude de (i). Levando em
conta (ii), vê-se então que (a') e (a'') são satisfeitas, logo, para
$n\geq\nu$, o homomorfismo natural $N_n\longrightarrow f_*(\mathcal G(n))$ é
um homomorfismo sobrejetivo entre Módulos localmente livres do mesmo posto e,
portanto, um isomorfismo; assim, tem-se
\[
f_*(\mathcal K(\nu+k))=\mathcal S_kR_\nu
\quad\text{para todo $k\geq0$},
\]
o que prova que $\mathcal G\in A_\nu(S)$ e que $R_\nu$ é o elemento de
$\operatorname{Grass}_{P(\nu)}(M_\nu)$ definido por $\mathcal G$.

Esse critério (i)--(ii) aplica-se igualmente à situação obtida após uma
mudança de base $S'\longrightarrow S$. Provaremos primeiro que o fato de a
condição (i) ser satisfeita após a mudança de base $S'\longrightarrow S$ se
exprime pela condição de que $S'\longrightarrow S$ se fatorize por certo
subpré-esquema $Z$ de $S$; obtido esse resultado, ficamos reduzidos
(substituindo $S$ por $Z$) ao caso em que a condição (i) já é satisfeita sobre
$S$ e, como ela é estável por mudança de base, resta exprimir a condição (ii).
Mas então, se $U$ denota o conjunto dos $s\in S$ tais que a cohomologia dos
feixes induzidos na fibra $X_s$ por $\mathcal G(n)$ e $\mathcal K(n)$ seja nula
em dimensão $>0$ para $n\geq\nu$, já se observou que $U$ é aberto, e a condição
(ii) será satisfeita após uma mudança de base $S'\longrightarrow S$ se e somente
se $S'\longrightarrow S$ se fatoriza por $U$, o que prova o Lema 3.3. Resta,
portanto, provar o lema seguinte:

\result{LEMA}{3.4} Sejam $S$ um pré-esquema localmente noetheriano, munido de
uma álgebra graduada quase-coerente $\mathcal S_*$ em graus positivos, gerada
por $\mathcal S_1$, $M_*$ um $\mathcal S_*$-Módulo graduado quase-coerente de tipo
finito, $P$ um polinômio com coeficientes racionais e $\nu$ um inteiro. Então
existe um subpré-esquema $Z$ de $S$ (evidentemente único)
\src{263}
com a seguinte propriedade: para todo pré-esquema $S'$ sobre $S$, para que
$M_n\otimes_{\mathcal O_S}\mathcal O_{S'}$ seja localmente livre de posto
$P(n)$ para todo $n\geq\nu$, é necessário e suficiente que
$S'\longrightarrow S$ se fatorize por $Z$.

Pode-se, evidentemente, supor $S$ afim, logo noetheriano. Então:

\result{LEMA}{3.5} Para todo inteiro $N\geq\nu$, seja $U_N$ o aberto de $S$
formado pelos $s\in S$ tais que
\[
\operatorname{rank}_{k(s)}
M_{ns}\otimes_{\mathcal O_{S,s}}k(s)\leq P(n)
\quad\text{para todo $\nu\leq n\leq N$}.
\]
Então a sequência decrescente dos abertos $U_N$ é estacionária.

Sabe-se [3], IV, que $S$ admite uma partição finita em subpré-esquemas
reduzidos $S_i$, tais que cada
$M\otimes_{\mathcal O_S}\mathcal O_{S_i}$ seja plano sobre $S_i$. Pode-se,
portanto, supor $M$ plano, logo os $M_n$ planos. Enfim, pode-se evidentemente
supor $S$ conexo. Mas então existem um inteiro $n_0$ e um polinômio $Q$ tais
que
\[
\operatorname{rank}_{k(s)}
M_{ns}\otimes_{\mathcal O_{S,s}}k(s)=Q(n)
\quad\text{para $n\geq n_0$},
\]
[5], III, \S~7. Suponhamos primeiro $P\neq Q$; então $P(n)\neq Q(n)$ para $n$
grande e, evidentemente, $U_N=\varnothing$ para $N$ grande; a fortiori, a
sequência dos $U_N$ é estacionária. Se, ao contrário, $P=Q$, então $U_N=U_{n_0}$ para $N\geq n_0$, e
a sequência dos $U_N$ é ainda estacionária.

Em particular, o conjunto $U_\infty$ dos $s\in S$ tais que
\[
\tag{*}
\operatorname{rank}_{k(s)}
M_{ns}\otimes_{\mathcal O_{S,s}}k(s)\leq P(n)
\quad\text{para todo $n\geq\nu$},
\]
sendo a interseção dos $U_N$, é aberto. Pode-se então, para a demonstração de
(3.4), substituir $S$ pelo aberto $U$, o que nos reduz ao caso em que a
desigualdade (*) é satisfeita em todo $s\in S$. Além disso:

\result{LEMA}{3.6} Seja $M$ um Módulo sobre um pré-esquema localmente
noetheriano $S$, e seja $r$ um inteiro. Então existe um subpré-esquema $Z$ de
$S$ (evidentemente único) com a seguinte propriedade: para todo $S'$ sobre
$S$, para que $M\otimes_{\mathcal O_S}\mathcal O_{S'}$ seja localmente
livre de posto $r$, é necessário e suficiente que $S'\longrightarrow S$ se
fatorize por $Z$. Se
\[
\operatorname{rank}_{k(s)}
M_s\otimes_{\mathcal O_{S,s}}k(s)\leq r
\quad\text{para todo $s$},
\]
então $Z$ é um subpré-esquema fechado de $S$ (supondo-se $M$ coerente).
\src{264}
Com efeito, o raciocínio precedente nos reduz ao caso em que a desigualdade
(*) vale para todo $s\in S$ (substituindo, se necessário, $S$ pela parte
aberta formada pelos $s$ em que a desigualdade é verificada). Pode-se então
supor que $M$ se insira em uma sequência exata
\[
\mathcal O_S^q\longrightarrow\mathcal O_S^r\longrightarrow M\longrightarrow0,
\]
e a condição considerada sobre os $S'$ sobre $S$ significa também que, na
sequência exata correspondente
$\mathcal O_{S'}^q\longrightarrow\mathcal O_{S'}^r
\longrightarrow M'\longrightarrow0$, a segunda seta é um isomorfismo, isto
é, a primeira é nula. Vê-se então que o subpré-esquema fechado $Z$ de $S$
definido pelo Ideal gerado pelos coeficientes da matriz que define o
homomorfismo $\mathcal O_S^q\longrightarrow\mathcal O_S^r$ satisfaz a
condição desejada.

Retomando então a demonstração de (3.4) no ponto em que a deixamos, denote-se
por $Z_n$ o subesquema fechado de $S$ associado, em virtude de (3.6), ao
Módulo $M_n$ e ao inteiro $r=P(n)$, e por $Z'_N$ o
$\operatorname{Inf}$ dos $Z_n$ para $\nu\leq n\leq N$. Então os $Z'_N$
formam uma sequência decrescente de subesquemas fechados de $Z$, portanto
necessariamente estacionária. Seja $Z$ o valor constante dos $Z'_N$ para $N$
grande. Esse é o $Z$ procurado em (3.4). Isso conclui a demonstração de (3.4)
e, portanto, também a de (3.3).

Demonstramos, pois, que
$\underline{\operatorname{Quot}}^P_{\mathcal F/X/S}$ é representável por um
$S$-pré-esquema $Q$ que é uma união crescente de subpré-esquemas abertos
$Q_\nu$ quase-projetivos sobre $S$. Para ir mais longe, é preciso invocar o
teorema 2.1, do qual se conclui facilmente que $Q$ é quase-compacto (pois é a
imagem do pré-esquema $S'$ de tipo finito sobre $S$ que parametriza a família
dos feixes quocientes dos $\mathcal F_K$ com polinômio de Hilbert $P$).
Portanto, $Q$ é igual a um dos $Q_\nu$, logo quase-projetivo sobre $S$. Para
demonstrar que é projetivo sobre $S$, resta demonstrar que é próprio sobre
$S$; para isso, basta invocar o critério valuativo de properidade sob a forma
[5], II, 7.3.8. Basta verificar o seguinte:

\result{LEMA}{3.7} Sejam $S$ o espectro de um anel de valorização discreta,
$s$ seu ponto genérico, $X$ um pré-esquema sobre $S$, $\mathcal F$ um Módulo
quase-coerente sobre $X$ e $\mathcal G_s$ um Módulo quociente quase-coerente
de $\mathcal F_s=\mathcal F\otimes_{\mathcal O_S}k(s)$ sobre $X_s$. Então
existe um único Módulo quociente quase-coerente $\mathcal G$ de $\mathcal F$,
plano sobre $S$, cuja restrição a $X_s$ seja $\mathcal G_s$.
\src{265}
Com efeito, se $\mathcal G_s=\mathcal F_s/\mathcal K_s$, basta considerar o
maior subfeixe $\mathcal K$ de $\mathcal F$ que induz $\mathcal K_s$, [5],
I, 9.4.2, e tomar $\mathcal G=\mathcal F/\mathcal K$. Verifica-se facilmente
que esse feixe convém.

O teorema 3.2 e, por conseguinte, (3.1) estão completamente demonstrados.

A demonstração prova, ao mesmo tempo, o seguinte:

\result{PROPOSIÇÃO}{3.8} Nas condições de (3.2), sejam
$Q=\operatorname{Quot}^P_{\mathcal F/X/S}$, $X_Q=X\times_SQ$,
$\mathcal F_Q=\mathcal F\otimes_{\mathcal O_S}\mathcal O_Q$, e seja
$\mathcal G$ o quociente coerente de $\mathcal F_Q$, plano sobre $Q$, com
polinômio de Hilbert relativo $P$, tal que $(Q,\mathcal G)$ represente o
funtor $\underline{\operatorname{Quot}}^P_{\mathcal F/X/S}$. Existe um
inteiro $\nu$ tal que, para $n\geq\nu$,
$(f_Q)_*(\mathcal G(n))$ seja um Módulo sobre $Q$ localmente livre de posto
$P(n)$ e muito amplo relativamente a $S$, isto é, que defina uma imersão de
$Q$ em um esquema grassmanniano $\operatorname{Grass}_{P(n)}(M)$ sobre $S$.
A fortiori, para $n\geq\nu$, o feixe
\[
\bigwedge^{P(n)}(f_Q)_*(\mathcal G(n))
\]
sobre $Q$ é invertível e muito amplo relativamente a $S$.

Pode-se, com efeito, reduzir, como em (3.2), ao caso em que $\mathcal F$ seja
plano sobre $S$; então basta tomar um inteiro $\nu$ tal que $A_\nu=A$ (com as
notações precedentes).

O caso de aplicação mais importante de (3.2) é aquele em que
$\mathcal F=\mathcal O_X$. Escreve-se então
\[
\operatorname{Quot}_{\mathcal O_X/X/S}=\operatorname{Hilb}_{X/S},
\qquad
\operatorname{Quot}^P_{\mathcal O_X/X/S}=\operatorname{Hilb}^P_{X/S},
\]
logo se tem uma decomposição
\[
\operatorname{Hilb}_{X/S}=\coprod_P\operatorname{Hilb}^P_{X/S}.
\]
Por definição, $\operatorname{Hilb}_{X/S}$ representa o funtor
$\underline{\operatorname{Hilb}}_{X/S}(S')$, conjunto dos subpré-esquemas
fechados de $X'=X\times_SS'$ que são planos sobre $S'$, e
$\operatorname{Hilb}^P_{X/S}$ representa o subfuntor correspondente aos
subpré-esquemas fechados que admitem um polinômio de Hilbert dado $P$. Esses
pré-esquemas são também chamados, respectivamente, \emph{pré-esquema de
Hilbert} de $X$ sobre $S$ e \emph{pré-esquema de Hilbert de índice $P$}. A
terminologia é justificada pelo papel desempenhado na teoria pelos polinômios
de Hilbert. Sua diferença de natureza em relação às variedades de Chow
clássicas (destinadas a parametrizar ciclos, e não variedades) é da mesma ordem
que aquela entre
\src{266}
o anel de Chow das classes de ciclos de uma variedade e o anel das classes
de feixes da variedade (tal como aparece no teorema de Riemann--Roch [1]);
com efeito, observe-se que, quando $X=\mathbf P^r_S$, sendo $S$ o espectro
de um corpo, o conhecimento do polinômio de Hilbert de um Módulo coerente
$\mathcal F$ sobre $X$ equivale também ao conhecimento das classes de Chern
de $\mathcal F$, ou ainda ao da classe de $\mathcal F$ no anel das classes de
feixes coerentes sobre $X$.

\result{OBSERVAÇÕES}{3.9} Observe-se também que a construção dos
$\operatorname{Quot}_{\mathcal F/X/S}$ e
$\operatorname{Quot}^P_{\mathcal F/X/S}$ foi reduzida ao caso em que
$X=\mathbf P^r_S$, $\mathcal F=\mathcal O_X^N$, sendo $\mathcal O_X(1)$ o
feixe muito amplo usual; mais precisamente, os
$\operatorname{Quot}_{\mathcal F/X/S}$ gerais realizam-se como
subpré-esquemas fechados dos precedentes. Como a formação dos
$\operatorname{Quot}^P_{\mathcal F/X/S}$ é evidentemente compatível com as
mudanças de base $S'\longrightarrow S$, vê-se que se está reduzido ao caso em
que, além disso, $S=\operatorname{Spec}(\mathbf Z)$. Em última análise, fica-se,
portanto, reduzido a estudar os esquemas projetivos sobre $\mathbf Z$:
\[
\mathfrak Q^P_{r,N}=
\operatorname{Quot}^{P}_{(\mathcal O_{\mathbf P^r_{\mathbf Z}})^N/
\mathbf P^r_{\mathbf Z}/\operatorname{Spec}(\mathbf Z)},
\]
e, mais particularmente, os esquemas de Hilbert absolutos:
\[
\operatorname{Hilb}^P_r=\mathfrak Q^P_{r,1}.
\]
Um estudo mais detalhado desses esquemas, começando pela determinação de suas
componentes conexas (são conexos?), e de suas componentes irredutíveis
(em virtude de Serre [7], pode haver componentes irredutíveis situadas
inteiramente sobre um número primo $p\neq0$), seria muito interessante.
Recordemos a questão de Weil sobre se as componentes irredutíveis das fibras de
$\operatorname{Hilb}^P_r$ sobre os $s\in\operatorname{Spec}(\mathbf Z)$
correspondem a extensões ``regulares'' do corpo primo, isto é, se são
``relativamente conexas''. Não se exclui que essas questões sejam mais
acessíveis para os esquemas de Hilbert do que para as ``variedades de Chow''.

\fsection{4}{Variantes}

\noindent a. Nas condições de (3.1), seja $U$ um aberto de $X$, e designemos
por $A'$ o subfuntor de
$A=\underline{\operatorname{Quot}}_{\mathcal F/X/S}$, tal que $A'(S')$ seja
o conjunto dos Módulos quocientes $\mathcal G$ de $\mathcal F'$, planos sobre
$S'$, cujo suporte está contido em $U'$. Vê-se imediatamente que $A'$ é
representável por uma parte aberta do pré-esquema
$\operatorname{Quot}_{\mathcal F/X/S}$ que representa $A$. Segue-se que os
teoremas (3.1) e (3.2)
\src{267}
continuam válidos supondo-se que $X$ seja quase-projetivo sobre $S$ em vez de
projetivo sobre $S$, quando se substitui também, nas conclusões, a palavra
``projetivo'' por ``quase-projetivo'', e se designa agora por
$\underline{\operatorname{Quot}}_{\mathcal F/X/S}(S')$ o conjunto dos
quocientes coerentes $\mathcal G$ de $\mathcal F'$, planos sobre $S'$, cujo
suporte é próprio sobre $S'$.

\medskip
\noindent b. De modo geral, podem-se impor aos quocientes $\mathcal G$ de
$\mathcal F'=\mathcal F\otimes_{\mathcal O_S}\mathcal O_{S'}$, planos sobre
$S'$, toda espécie de condições naturais adicionais, estáveis por mudança de
base, obtendo-se assim tantos subfuntores de
$\underline{\operatorname{Quot}}_{\mathcal F/X/S}$ quantos se pretende
representar. O critério usual permite, em muitos casos, demonstrar que ainda se
obtêm funtores representáveis por partes abertas de
$\operatorname{Quot}_{\mathcal F/X/S}$. Isso ocorre, em particular, quando se
impõe uma das seguintes propriedades adicionais:

\smallskip
\noindent $1^\circ$ As dimensões dos ciclos primos associados aos Módulos
$\mathcal G_{s'}$ ($s'\in S'$) induzidos nas fibras $X'_{s'}$ pertencem a um
conjunto dado de inteiros.

\smallskip
\noindent $2^\circ$ (Quando $\mathcal F=\mathcal O_X$, portanto $\mathcal G$
corresponde a um subpré-esquema fechado $Y$ de $X'$); $Y$ é um pré-esquema
suave ([3], IV) sobre $S$, resp. normal sobre $S$ (isto é, as fibras $Y_{s'}$
são normais ``sobre $k(s')$'', isto é, são normais após toda extensão do corpo
de base), resp. (quando $X$ é plano sobre $S$) é uma $S$-interseção completa
local em $X$ relativamente a $S$ (isto é, as fibras $Y_{s'}$ são interseções
completas locais nos $X_{s'}$).

Outras condições fariam intervir propriedades de natureza cohomológica nos
Módulos $\mathcal G_{s'}$ induzidos nos $X'_{s'}$, etc. Naturalmente, a
conjunção de condições, cada uma das quais representada por um aberto $U_i$ de
$\operatorname{Quot}_{\mathcal F/X/S}$, é representada pelo aberto que é sua
interseção. Por exemplo, considerando, para todo $S'$ sobre $S$, o conjunto dos
subpré-esquemas fechados $Y$ de $X'=X\times_SS'$ que são recobrimentos étales
([3], I) de posto dado $r$ de $S'$, obtém-se um funtor contravariante
representável em $S'$.

\medskip
\noindent c. Os pré-esquemas $\operatorname{Hom}_S(X,Y)$,
$\Pi_{X/S}(Z/X)$, $\operatorname{Isom}_S(X,Y)$, definidos em [2], II, C,
n\textsuperscript{o} 2, existem sob hipóteses de projetividade convenientes e
realizam-se como partes abertas de pré-esquemas de Hilbert convenientes. Como se
tem
\[
\operatorname{Hom}_S(X,Y)=\Pi_{X/S}((X\times_SY)/X),
\]
o caso de $\operatorname{Hom}_S(X,Y)$ reduz-se ao de $\Pi_{X/S}(Z/X)$.
Observa-se então que, para todo $S'$ sobre $S$, o conjunto das seções de
$Z'=Z\times_SS'$ sobre $X'=X\times_SS'$ está em correspondência biunívoca com
o conjunto dos subpré-esquemas $\Gamma$ de $Z'$ (necessariamente fechados se
$Z$ é separado sobre $X$) tais que o morfismo $\Gamma\longrightarrow X'$
induzido por $Z'\longrightarrow X'$ seja um isomorfismo. Dessa maneira,
\src{268}
quando $X$ é plano e próprio sobre $S$ e $Z$ é quase-projetivo sobre $S$,
$\Pi_{X/S}(Z/X)$ existe e realiza-se como um subpré-esquema aberto de
$\operatorname{Hilb}_{Z/S}$. Portanto, quando $X$ é projetivo e plano sobre
$S$ e $Y$ é quase-projetivo sobre $S$, $\operatorname{Hom}_S(X,Y)$ existe e
realiza-se como um subpré-esquema aberto de
$\operatorname{Hilb}_{(X\times_SY)/S}$. Quando $X$ e $Y$ são ambos projetivos
sobre $S$, segue-se imediatamente que $\operatorname{Isom}_S(X,Y)$ também
existe e é representado por uma parte aberta de
$\operatorname{Hom}_S(X,Y)$. Do mesmo modo, quando $X$ é plano e projetivo
sobre $S$ e $Y$ é quase-projetivo sobre $S$, o $S$-pré-esquema
$\operatorname{Imm}_S(X/Y)$, correspondente ao subfuntor do funtor representado
por $\operatorname{Hom}_S(X,Y)$ que corresponde aos $S'$-homomorfismos
$X'\longrightarrow Y'$ que são imersões, é também representável por uma
parte aberta de $\operatorname{Hom}_S(X,Y)$.

Seja $\mathcal L$ (resp. $\mathcal M$) um feixe invertível sobre $X$ (resp.
$Y$), muito amplo relativamente a $S$, donde um feixe muito amplo
$\mathcal L\otimes_{\mathcal O_S}\mathcal M$ sobre $X\times_SY$ relativamente a
$S$. Por conseguinte, para todo polinômio $P$ com coeficientes racionais,
$\operatorname{Hilb}^P_{(X\times_SY)/S}$ está definido e é um pré-esquema
quase-projetivo sobre $S$. Ele induz, portanto, sobre
$\operatorname{Hom}_S(X,Y)$ uma parte simultaneamente aberta e fechada,
quase-projetiva sobre $S$, que denotaremos por
$\operatorname{Hom}_S(X,Y)^P$. Assim, as seções de
$\operatorname{Hom}_S(X,Y)^P$ sobre $S$ são os $S$-morfismos
$g:X\longrightarrow Y$ tais que, para todo inteiro $n$, se tenha
\[
\chi\!\left(\left(\mathcal L\otimes_{\mathcal O_X}g^*(\mathcal M)\right)^{
\otimes n}\right)=P(n).
\]

Obtêm-se dessa maneira generalizações do teorema de Matsusaka, que afirma
que os automorfismos de uma variedade projetiva «polarizada» formam um grupo
algébrico, afirmação que assume aqui um sentido evidentemente mais preciso,
pois dispomos de uma definição desse grupo como solução de um problema
universal. Note-se, aliás, que, sobre um corpo algebricamente fechado, o grupo
dos automorfismos considerado anteriormente é o que se deduz do «verdadeiro»
definido aqui passando ao quociente pelos elementos nilpotentes; isso explica
por que há pouca chance de que as construções mais antigas possam ser feitas
sobre um corpo de base não perfeito, pois o Ideal dos elementos nilpotentes que
aparece após uma extensão do corpo de base não é necessariamente «definido
sobre $k$». Essa mesma observação aplica-se, aliás, também à maioria das construções
de estilo antigo.

\fsection{5}{Estudo diferencial dos esquemas de Hilbert}

Ele decorre do resultado seguinte:
\src{269}
\result{PROPOSIÇÃO}{5.1} Sejam $S$ um pré-esquema, $S_0$ um subpré-esquema
definido por um Ideal quase-coerente $\mathcal J$ de quadrado nulo, $X$ um
$S$-pré-esquema, $\mathcal F$ um Módulo quase-coerente sobre $X$,
$X_0=X\times_SS_0$ e
$\mathcal F_0=\mathcal F\otimes_{\mathcal O_S}\mathcal O_{S_0}$, e, por fim,
$\mathcal G_0=\mathcal F_0/\mathcal K_0$ um Módulo quociente quase-coerente de
$\mathcal F_0$, plano sobre $S_0$. Para todo aberto $U$ de $X$, seja
$\mathcal E(U)$ o conjunto dos Módulos quocientes quase-coerentes $\mathcal G$ de
$\mathcal F|U$, planos sobre $S$, e tais que
$\mathcal G\otimes_{\mathcal O_S}\mathcal O_{S_0}=\mathcal G_0$; assim, quando
$U$ varia, os $\mathcal E(U)$ são as seções de um feixe $\mathcal E$ sobre
$U$. Posto isso, o feixe de grupos
\[
\mathcal A=\operatorname{Hom}_{\mathcal O_{X_0}}
(\mathcal K_0,\mathcal G_0\otimes_{\mathcal O_{S_0}}\mathcal J)
\]
age naturalmente sobre $\mathcal E$, que se torna assim um feixe
«formalmente homogêneo principal sob $\mathcal A$» (isto é, para todo aberto
$U$ em $X$, $\mathcal E(U)$ é vazio ou um conjunto homogêneo principal sob
$\mathcal A(U)$).

Conclui-se daí:

\result{COROLÁRIO}{5.2} Suponhamos que exista localmente sobre $X$ um
prolongamento $\mathcal G$ de $\mathcal G_0$ a um quociente de $\mathcal F$
plano sobre $S$ (isto é, que as fibras do feixe $\mathcal E$ sejam não vazias).
Então existe uma classe de obstrução canônica
\[
c(\mathcal G_0)\in H^1(X,\mathcal A),
\]
cujo anulamento é necessário e suficiente para a existência de um
prolongamento global $\mathcal G$ de $\mathcal G_0$ a um quociente de
$\mathcal F$ plano sobre $S$. Se essa classe se anula, então o conjunto
$\mathcal E(X)$ de todos os prolongamentos possíveis é um conjunto homogêneo
principal sob
\[
\mathcal A(X)=\operatorname{Hom}_{\mathcal O_X}
(\mathcal K_0,\mathcal G_0\otimes_{\mathcal O_{S_0}}\mathcal J).
\]

A existência do prolongamento global fica então garantida, em particular, se
$H^1(X,\mathcal A)=0$.

\result{COROLÁRIO}{3.3} Suponhamos que
$Q=\operatorname{Quot}_{\mathcal F/X/S}$ exista (cf. 4 (a)) -- por exemplo,
que $X$ seja quase-projetivo sobre um $S$ localmente noetheriano e que
$\mathcal F$ seja coerente. Seja $x\in Q$ correspondente a uma extensão
residual $K=k(x)$ de $k(s)$ ($s\in S$):
\src{270}
portanto, $x$ é definido por um Módulo quociente coerente
$\mathcal G_0=\mathcal F_0/\mathcal K_0$ do Módulo
$\mathcal F_0=\mathcal F_K$ sobre o $K$-pré-esquema $X_K$. Seja $\mathcal A$
o feixe coerente sobre $X_K$ definido por
\[
\mathcal A=\operatorname{Hom}_{\mathcal O_{X_0}}(\mathcal K_0,\mathcal G_0).
\]
Então o espaço tangente de Zariski da fibra $Q_s$ no ponto $x$ (dual sobre
$K$ de $\mathfrak m/\mathfrak m^2$, onde $\mathfrak m$ é o ideal maximal de
$\mathcal O_{Q_K,x}$) é canonicamente isomorfo a $H^0(X_K,\mathcal A)$.

O resultado que fornece o espaço tangente de Zariski pode ser generalizado e
fornece uma caracterização, para um $S$-morfismo dado
$g:S'\longrightarrow Q$, isto é, uma seção $g'$ de $Q'=Q\times_SS'$ sobre
$S'$, do Módulo
\[
\Omega=g^*(\Omega^1_{Q/S})=g'^*(\mathcal J/\mathcal J^2)
\]
(onde $\mathcal J$ é o Ideal sobre $Q'$ definido pela seção $g'$ de $Q'$ sobre
$S'$), pela fórmula funtorial no Módulo coerente $\mathcal M$ sobre $S'$:
\[
\operatorname{Hom}_{\mathcal O_{S'}}(\Omega,\mathcal M)
\simeq H^0(X',\mathcal A),
\]
onde $\mathcal A$ é ainda o Módulo sobre $X'=X\times_SS'$ definido por
\[
\mathcal A=\operatorname{Hom}_{\mathcal O_{X'}}
(\mathcal K,\mathcal G\otimes_{\mathcal O_S}\mathcal M),
\]

($\mathcal G=\mathcal F'/\mathcal K$ sendo o Módulo quociente de
$\mathcal F'=\mathcal F\otimes_{\mathcal O_S}\mathcal O_{S'}$ que corresponde
a $g$). Com efeito, basta aplicar (5.1), substituindo ali $S_0$ por $S'$ e $S$
pelo pré-esquema $D(\mathcal M)=(S',\mathcal O_{S'}+\mathcal M)$, onde
$\mathcal M$ é considerado como um Ideal de quadrado nulo.

Quando, em (5.1), se tem $\mathcal F=\mathcal O_X$, portanto o dado de
$\mathcal G_0$ corresponde ao dado de um subpré-esquema fechado $Y_0$ de
$X_0$, plano sobre $S_0$, definido pelo Ideal
$\mathcal J_0=\mathcal K_0$, então (*) fornece
\[
\mathcal A=\operatorname{Hom}_{\mathcal O_{X_0}}
(\mathcal J_0/\mathcal J_0^2,
\mathcal O_{Y_0}\otimes_{\mathcal O_{S_0}}\mathcal J),
\]
onde $\mathcal J/\mathcal J^2$ se interpreta como o feixe conormal a $Y_0$ em
$X_0$, que se denota também por
\src{271}
$\check{\mathcal N}_{Y_0/X_0}$; é então conveniente considerar $\mathcal A$
como um Módulo sobre $Y_0$ e calcular $H^0$ e $H^1$ sobre $Y_0$. Ademais, se
$Y_0$ é localmente uma interseção completa em $X_0$, sendo $X$ plano sobre
$S$, então, em (5.1), fica garantida a possibilidade de prolongamento local;
por outro lado, $\mathcal J/\mathcal J^2$ é localmente livre sobre $Y_0$ e
pode-se escrever
\[
\mathcal A=\mathcal N_{Y_0/X_0}\otimes_{\mathcal O_{S_0}}\mathcal J,
\]
onde o primeiro fator do segundo membro é o feixe normal a $Y_0$ em $X_0$.
Utilizando o critério fundamental de lisura ([3], III, 3.1), obtém-se, por
exemplo:

\result{COROLÁRIO}{5.4} Nas condições de (5.3), suponhamos que
$\mathcal F=\mathcal O_X$, que $X$ seja plano sobre $S$ e que o
subpré-esquema fechado $Y_0$ de $X_0$ que corresponde a $\mathcal G_0$ seja
localmente uma interseção completa. Então o espaço tangente de Zariski a $Q_s$
em $x$ é canonicamente isomorfo a
$H^0(Y_0,\mathcal N_{Y_0/X_0})$. Se
$H^1(Y_0,\mathcal N_{Y_0/X_0})=0$, então o pré-esquema de Hilbert
$\operatorname{Hilb}_{X/S}$ é liso sobre $S$ no ponto $x$.
($\mathcal N_{Y_0/X_0}$ é o feixe normal a $Y_0$ em $X_0$.)

\result{OBSERVAÇÃO}{5.5} Este enunciado aplica-se, em particular, quando $Y_0$
é uma interseção completa em $X_0$ definida por uma equação, isto é, é um
``divisor de Cartier'' positivo. Então $\mathcal N_{Y_0/X_0}$ é isomorfo ao
feixe sobre $Y_0$ induzido pelo feixe invertível $\mathcal J^{-1}$ sobre $X_0$
definido pelo divisor $Y_0$. É essa a situação encontrada, em particular, no
estudo das famílias de divisores positivos sobre uma variedade projetiva não
singular $X_0$. O isomorfismo entre o espaço tangente de Zariski no ponto $x$
de $Q$ (ou, se se preferir, do aberto $D$ de $Q$ que corresponde aos divisores)
e $H^0(Y_0,\mathcal N_{Y_0/X_0})$ era conhecido em geometria algébrica
clássica sob o nome de ``homomorfismo característico'' do primeiro no segundo.
Ele só era definido quando $x$ era um ponto liso da variedade de parâmetros
$T$ de uma ``família contínua completa'' de divisores, isto é, de nosso ponto
de vista, de uma componente irredutível do esquema $D$, munida da estrutura
reduzida induzida. O espaço tangente a $T$ em $x$ é então um subespaço do
espaço tangente a $D$ em $x$; portanto, o homomorfismo característico clássico
é de fato injetivo, mas só é sobrejetivo sob condições adicionais, por exemplo,
se $D$ é integral em $x$. De fato, Zappa [8] construiu um exemplo
\src{272}
(com $X$ uma superfície projetiva não singular sobre o corpo dos complexos) no
qual, mesmo no ponto genérico de $T$, o homomorfismo característico não é
sobrejetivo. Isso significa, portanto, que $D$ não é integral nem sequer no
ponto genérico da componente irredutível considerada. Isso mostra de maneira
particularmente impressionante como as variedades com elementos nilpotentes
são necessárias para compreender fenômenos da mais clássica teoria das
superfícies.

\subsection*{5.6}

Em (5.4), deu-se um critério de lisura, aplicável em particular aos esquemas de
divisores. Kodaira deu em [6] um critério diferente, a saber, o anulamento de
$H^1(X_0,\mathcal L)$, onde $\mathcal L=\mathcal J_0^{-1}$ é o feixe
invertível sobre $X_0$ definido pelo divisor $Y_0$; critério válido quando $S$
é o espectro de um corpo de característica 0 e demonstrado em [6] por via
transcendente no caso em que o corpo de base é $\mathbf C$. Observemos aqui
que, em geral, sendo $S$ novamente arbitrário, a condição de Kodaira é
suficiente para que o morfismo canônico
\[
D\longrightarrow\operatorname{Pic}_{X/S}
\]
do pré-esquema dos divisores para o pré-esquema de Picard de $X/S$ seja liso
no ponto $x$ considerado (como se verifica facilmente pelo critério usual de
lisura, uma vez estabelecida a existência de $\operatorname{Pic}_{X/S}$). Se,
portanto, além disso, $\operatorname{Pic}_{X/S}$ é liso sobre $S$ no ponto
imagem de $x$ (por exemplo, se $\operatorname{Pic}_{X/S}$ é liso sobre $S$),
então $D$ é liso sobre $S$ em $x$. Por outro lado, Cartier demonstrou que todo
pré-esquema em grupos localmente de tipo finito sobre um corpo $k$ de
característica 0 é liso sobre $k$. Combinando esses dois resultados, recupera-se
o resultado de Kodaira. Observe-se que decorre dessas observações que, sobre um
corpo $K$ de característica $p>0$, quando $\operatorname{Pic}_{X/S}$ não é
liso sobre $k$ (o que ocorre quando $X$ é a superfície de Igusa), a condição
$H^1(X_0,\mathcal L)=0$ implica, ao contrário, que $D$ não é liso em $x$, e
nem sequer é reduzido em $x$ se $K$ é algebricamente fechado.

Para terminar, assinalemos ainda o resultado seguinte, que desempenha um papel
importante no estudo diferencial dos espaços fibrados:

\result{PROPOSIÇÃO}{5.7} Sejam $X$ um pré-esquema finito e plano sobre $S$
localmente noetheriano, e $Z$ um pré-esquema sobre $S$, tal que
$\Pi_{X/S}(Z/X)$ exista (o que ocorre se $Z$ é quase projetivo sobre $X$). Se,
então, $Z$ é liso sobre $X$, $\Pi_{X/S}(Z/X)$ é liso sobre $S$.

É uma consequência imediata da definição e do critério usual de lisura ([3],
III, 3.1). Observemos que, quando $X$ é finito e plano sobre $S$, a questão da
existência de $Z/K$ pode ser tratada de maneira muito elementar, sem
\src{273}
utilizar a teoria dos esquemas de Hilbert. Verifica-se, por exemplo, que, se
$X$ é radicial sobre $S$, então $\Pi_{X/S}(Z/X)$ existe sem nenhuma restrição
sobre $Z$. Por exemplo, seja $T$ um $S$-pré-esquema, e seja $T_n$ a
``vizinhança infinitesimal de ordem $n$'' da diagonal de $T\times_ST$ em
$T\times_ST$, munida dos morfismos $p_1,p_2:T_n\longrightarrow T$ induzidos
pelas duas projeções. Considera-se $T_n$ como um pré-esquema finito sobre $T$
por meio de $p_1$; suponhamo-lo, além disso, plano sobre $T$ (o que ocorre se
$T$ é liso sobre $S$). Para todo pré-esquema $X$ sobre $T$, ponhamos
\[
(X/T/S)^{(n)}=
\Pi_{T_n/S}\bigl(p_2^*(X/T)/T_n\bigr),
\]

este é um pré-esquema sobre $T$ chamado \emph{fibrado dos germes de seções de
ordem $n$ de $X$ sobre $T$} (relativamente a $S$). Ele depende funtorialmente
de $X$ e é liso sobre $T$ se $X$ o é.

\fsection{6}{Relações com a noção de norma e os produtos simétricos}

Sejam $S$ um pré-esquema, $X$ e $Y$ $S$-pré-esquemas,
\[
u:(X/S)^n\longrightarrow Y
\]
um $S$-morfismo simétrico da $n$-ésima potência cartesiana de $X/S$ em $Y$.
Suporemos, para simplificar, $S$ localmente noetheriano, e $X$ e $Y$ de tipo
finito sobre $S$. Pode-se então, a todo Módulo coerente $\mathcal F$ sobre $X$,
com suporte finito sobre $S$, plano sobre $S$ e de posto relativo sobre $S$
igual a $n$ (isto é, tal que $f_*(\mathcal F)$ seja um Módulo localmente livre
de posto $n$ sobre $S$), fazer corresponder de maneira natural uma seção de $Y$
sobre $S$:
\[
\pi^u_{X/S}(\mathcal F)\in\Gamma(Y/S).
\]
Não daremos aqui a definição formal, limitando-nos a assinalar que o formalismo
ao qual se chega é uma generalização natural do formalismo usual das normas e
dos traços. Quando a $n$-ésima potência simétrica de $X$ sobre $S$ existe (por
exemplo, quando as órbitas do grupo simétrico $\mathfrak S_n$ que opera sobre
$(X/S)^n$ estão contidas em abertos afins), pode-se tomar para $Y$ essa
potência simétrica $\operatorname{Symm}^n_S(X)$, e obtém-se um elemento canônico
\[
\pi_{X/S}(\mathcal F)\in\Gamma(\operatorname{Symm}^n_S(X)/S),
\]
\src{274}
que permite recuperar os $\pi^u_{X/S}(\mathcal F)$. Outro caso importante é
aquele em que $X$ é um monoide comutativo sobre $S$, e $X=Y$, provindo o
morfismo $u$ da lei de composição de $X$. Escrever-se-á então simplesmente
$\pi(\mathcal F)$ para a seção de $X$ sobre $S$ associada ao Módulo
$\mathcal F$ sobre $X$.

Suponhamos agora que se tenha um Módulo coerente $\mathcal F$ sobre $X$, tal
que $\operatorname{Quot}_{\mathcal F/X/S}$ exista, ou pelo menos tal que o
funtor $\underline{\operatorname{Quot}}^n_{\mathcal F/X/S}$, que associa a
todo $S'$ sobre $S$ o conjunto dos feixes quocientes coerentes $\mathcal M$ de
$\mathcal F'=\mathcal F\otimes_{\mathcal O_S}\mathcal O_{S'}$ que são planos
sobre $S'$ e de posto relativo $n$, seja representável por um $S$-pré-esquema
$\operatorname{Quot}^n_{\mathcal F/X/S}$. (Quando $X$ é quase projetivo sobre
$S$, então $\operatorname{Quot}^n_{\mathcal F/X/S}$ existe de fato e não é
outro, com as notações do número 3, senão
$\operatorname{Quot}^P_{\mathcal F/X/S}$, onde $P$ é o polinômio reduzido ao
termo constante $n$.) Como a formação dos $\pi^u_{X/S}(\mathcal M)$ é
compatível com a mudança de base, obtém-se então um morfismo canônico
\[
\pi^u_{X/S}:\operatorname{Quot}^n_{\mathcal F/X/S}\longrightarrow Y,
\]
em particular, se a $n$-ésima potência simétrica de $X$ sobre $S$ existe:
\[
\pi_{X/S}:\operatorname{Quot}^n_{\mathcal F/X/S}
\longrightarrow\operatorname{Symm}^n_S(X).
\]
O caso mais importante é aquele em que $\mathcal F=\mathcal O_X$, que fornece
um morfismo:
\[
\pi_{X/S}:\operatorname{Hilb}^n_{X/S}
\longrightarrow\operatorname{Symm}^n_S(X).
\]
Este é evidentemente um isomorfismo para $n=0$ ou $n=1$. Mas, para
$n\geq1$, mesmo se $S$ é o espectro
de um corpo $k$, e quando $X$ é liso sobre $S$, ele não é, em geral, um
isomorfismo, nem sequer um morfismo injetivo, pois um subesquema de dimensão 0
de $X$ (correspondente, por exemplo, a um ideal $\mathfrak i$ primário para o
ideal maximal em um anel local $\mathcal O_{X,x}$, para um ponto fechado $x$ de
$X$) não é conhecido quando se conhece o ciclo que ele define (no caso, quando
se conhece a codimensão sobre $k$ de $\mathfrak i$ em $\mathcal O_{X,x}$).
Pode-se dizer apenas o seguinte (sendo $S$ novamente arbitrário):

\medskip
\noindent a. Se $X$ é liso sobre $S$, então o morfismo-norma define um
isomorfismo do aberto de $\operatorname{Hilb}^n_{X/S}$ que corresponde à
classificação dos recobrimentos étales de posto $n$ contidos em $X$ (cf.
n\textsuperscript{o} 4, (b)) com o aberto de
$\operatorname{Symm}^n_S(X)$ que corresponde aos $n$-ciclos sem componentes
múltiplas.
\src{275}

\medskip
\noindent b. Se, além disso, $X$ tem dimensão relativa 1 sobre $S$, então o
morfismo norma define até mesmo um isomorfismo de
$\operatorname{Hilb}^n_{X/S}$ com $\operatorname{Symm}^n_{X/S}$.

Este último fato deve-se a que um subesquema de dimensão 0 de uma curva
algébrica não singular é conhecido quando se conhece o ciclo correspondente.
A mesma observação aplica-se, aliás, mais geralmente aos divisores de Cartier
positivos sobre um esquema algébrico não singular (e não se exclui que, nesse
caso muito particular, a variedade de Chow forneça a mesma coisa que o esquema
de Hilbert).

\fsection{7}{Complementos e questões}

Como observou J.-P. Serre, resulta de um exemplo bem conhecido de Nagata que se
pode encontrar um esquema $S$, espectro de um corpo $k$, um $S$-esquema $S'$,
espectro de uma extensão quadrática $k'$ de $k$, enfim, um $S'$-esquema $X$
próprio e liso (mas não projetivo) de dimensão 3, tal que
$\Pi_{S'/S}(X/S')$ não existe. Isso implica, a fortiori, que o esquema de
Hilbert $\operatorname{Hilb}^2_{X/S}$ não existe (nem sequer o $k$-esquema que
representaria os recobrimentos étales de posto 2 de $S$ contidos em $X$, nem,
a fortiori, o quadrado simétrico de $X$, cf. número precedente). Isso impõe,
portanto, limitações sérias às possibilidades de construções não projetivas em
geometria algébrica. (É, contudo, plausível que tais limitações não se
apresentem em geometria analítica, assim como não se apresentam em geometria
formal (cf. [2], II).) Em contrapartida, se $X$ é um esquema próprio sobre o
espectro $S$ de um corpo $k$, e se $Z$ é quase projetivo sobre $X$, então
$\Pi_{X/S}(Z/X)$ existe e é um esquema, soma disjunta de uma sequência de
esquemas quase projetivos sobre $S$ (como no caso projetivo (3.1)). Para vê-lo,
reduz-se, com efeito, ao caso em que o próprio $X$ é projetivo, dominando $X$
por um $S$-esquema projetivo $X'$; não daremos aqui os detalhes da demonstração,
que utiliza também o resultado de fatoração de um morfismo finito assinalado em
[2], I, A, n\textsuperscript{o} 2 (b). O êxito do método deve-se ao fato de que,
sendo $S$ o espectro de um corpo, o $X'$ que intervém no lema de Chow será
automaticamente plano sobre $S$.

Ignoro se o resultado continua válido sem hipótese sobre $S$, supondo apenas
$X$ próprio e plano sobre $S$, e $Z$ quase projetivo sobre $X$. Um caso
importante nas aplicações é aquele em que $Z$ é um subesquema fechado de $X$;
se, então, $\Pi_{X/S}(Z/X)$ existe, ele é necessariamente um subesquema fechado
de $S$. Pode-se
\src{276}
construí-lo diretamente de maneira bastante simples quando $X$ é projetivo
sobre $S$, sem utilizar a teoria dos esquemas de Hilbert, e o método empregado
mostra, mais geralmente, que, se $Z$ é afim sobre $X$, então
$\Pi_{X/S}(Z/X)$ existe e é afim sobre $S$. Ele mostra igualmente que, se $X$
é próprio e plano sobre $S$ (não necessariamente projetivo sobre $S$), então,
para todo fibrado vetorial $Z$ localmente trivial sobre $X$,
$\Pi_{X/S}(Z/X)$ existe e é um fibrado vetorial sobre $S$. Seria desejável que
esses resultados fossem retomados e unificados.

\section*{BIBLIOGRAFIA}
\addcontentsline{toc}{section}{Bibliografia}

\begin{description}
\item[{[1]}] BOREL (A.) e SERRE (J.-P.). -- Le théorème de Riemann--Roch,
\emph{Bull. Soc. math. France}, v. 86, 1958, p. 97--136.

\item[{[2]}] GROTHENDIECK (Alexander). -- Technique de construction et
théorèmes d'existence en géométrie algébrique, I, II, III, \emph{Séminaire
Bourbaki}, v. 13, 1959/60, n\textsuperscript{o} 190, 29 p.,
n\textsuperscript{o} 195, 22 p.; e v. 14, 1960/61,
n\textsuperscript{o} 212, 20 p.

\item[{[3]}] GROTHENDIECK (Alexander). -- Séminaire de Géométrie algébrique,
I, II, III, IV. -- Paris, Institut des hautes Études scientifiques, 1960/61
(mimeografado).

\item[{[4]}] GROTHENDIECK (Alexander). -- Technique de construction en
géométrie analytique, IV, V, \emph{Séminaire Cartan}, v. 13, 1960/61,
n\textsuperscript{o} 11 e 12.

\item[{[5]}] GROTHENDIECK (A.) e DIEUDONNÉ (J.). -- \emph{Éléments de
géométrie algébrique}, I: Le langage des schémas. -- Paris, Presses
universitaires de France, 1960 (Institut des hautes Études scientifiques,
Publications mathématiques, 4); II, III, IV (a aparecer).

\item[{[6]}] KODAIRA (K.). -- Characteristic linear systems of complete
continuous systems, \emph{Amer. J. of Math.}, v. 78, 1956, p. 716--744.

\item[{[7]}] SERRE (Jean-Pierre). -- Exemples de variétés projectives en
caractéristique $p$ non relevables en caractéristique 0, \emph{Proc. Nat. Acad.
Sc. U.S.A.}, v. 47, 1961, p. 108--109.

\item[{[8]}] ZAPPA (G.). -- Sull'esistenza, sopra le superficie algebriche, di
sistemi continui completi infiniti, la cui curva è a serie caratteristica
incompleta, \emph{Pont. Acad. Sc. Acta}, v. 9, 1945, p. 91--93.
\end{description}

\section*{ADENDO}
\addcontentsline{toc}{section}{Adendo}
\begin{center}
{\small [acrescentado durante a correção das provas]}
\end{center}

Parece agora que as conjecturas de [2], III, n\textsuperscript{o} 8, são falsas,
mesmo para variedades não singulares sobre um corpo de característica 0, tanto
no que concerne à existência quanto à quase projetividade do quociente, e mesmo
quando $G$ opera com um grafo fechado.
